The continuous triangular Hilbert form has the same cyclic arrangement
of inputs as \eqref{eq:local-form}, integrated over \(\R^3\) against the
principal-value kernel \(1/(x_0+x_1+x_2)\). Its symmetric \(L^3\) estimate is formulated by
Thiele in \cite[Section~9, Problem~13]{ThieleProblem2017}. That
contribution also displays the precise dyadic model
\eqref{eq:given-form}, with one coefficient for each scale, and asks
for the analogous bound. Theorem~\ref{thm:main} answers this dyadic
question affirmatively and gives the stronger sum of absolute local
contributions. The continuous form and its dyadic model are distinct
operators; the theorem here is a statement about the latter.

Kova\v{c}, Thiele and Zorin-Kranich introduced a dyadic triangular
model with coefficients allowed to vary among interval triples
\cite[Equations~(1.5)--(1.6)]{KTZ2015}. Their Theorem~1.7 gives
uniform estimates when one input has either a specified dependence
on a Walsh linear combination of its coordinates or a specified
fiberwise character structure. Those cases already encompass dyadic
versions of Carleson's operator and uniform bounds for the bilinear
Hilbert transform, as explained in their Appendix~A. Our estimate
allows all three inputs to vary freely. The cyclic operator-trace
representation used below also appears in their Equation~(1.5).

Another line of work measures cancellation through the growth of
truncated norms. Tao obtained sublogarithmic growth for
one-dimensional multilinear Hilbert transforms
\cite[Theorem~1.2]{Tao2016}; Zorin-Kranich proved the corresponding
sublinear growth in the number of smooth scales for simplex Hilbert
transforms \cite[Theorem~1.3]{ZK2017}. Durcik, Kova\v{c} and Thiele
then obtained power-type cancellation in both the continuous and
dyadic settings \cite[Theorems~1 and~4]{DKT2019}. For the triangular
dyadic form, their bound over \(m\) consecutive scales has a factor
\(m^{1/2}\) at exponents \((4,4,2)\); cyclic permutation and multilinear
interpolation give the same growth at \((3,3,3)\). Theorem~\ref{thm:main}
removes this scale dependence at the symmetric tuple.

The strategy of controlling an entangled form by the change of a
nonnegative energy has precedents in Kova\v{c}'s Bellman-function
framework \cite[Section~2]{Kovac2011} and his telescoping identity
for the twisted paraproduct \cite[Lemma~2]{Kovac2012}.
These provide a methodological precedent, rather than a bound for the
present triangular form. In particular, Kova\v{c} explicitly identifies
the obstruction to applying that framework to non-bipartite graphs,
including the triangle, in \cite[Section~6.2]{Kovac2011}.
Here the energy is a trace of a matrix power, and the averaging is
over the four children selected by the XOR constraint.

The matrix energy and the dimension-independent trace inequality are
taken from the companion manuscript \cite[Sections~2--3, Lemmas~2.1,
2.2 and~3.1]{OpenAI2026}; we reproduce their real finite-dimensional
proofs. Their underlying matrix calculus is classical: the
square-root derivative is a special case of the divided-difference
formula in \cite[Section~1, Equation~(1)]{Daletskii1957}, and the order
comparison uses monotonicity of the positive square root
\cite[Theorem, p.~309]{Pedersen1972}.
